% =============================================================================
% Project : THE PHILOSOPHY OF VIABILITY
% Book    : Book I — Ontology of Viability
% Book ID : VS.B01
% Part    : Part I — Identity and Manifestation
% Part ID : VS.B01.P01
% Chapter : Chapter 08 — Comparative Philosophical Examination
% Chapter ID : VS.B01.P01.C08
%
% File :
%   2_MAINMATTER/PART_I_IDENTITY_AND_MANIFESTATION/
%   CHAPTER_08_COMPARATIVE_PHILOSOPHICAL_EXAMINATION/
%   sections/section_11_heidegger.tex
%
% Section :
%   Section 11 — Martin Heidegger
%
% Structural Identifier :
%   VS.B01.P01.C08.S11
%
% Purpose :
%   To reconstruct Heidegger's question of Being, ontological
%   difference, analytic of Dasein, being-in-the-world, worldhood,
%   equipmental involvement, being-with, care, authenticity,
%   being-toward-death, temporality, disclosure, and later
%   thinking; and to assess their significance and limits for
%   the ontology of Being, Identity, Determination, and
%   Manifestation developed in Book I.
%
% Method :
%   Historically responsible reconstruction, primary-text
%   examination, interpretive qualification, symmetrical
%   dialectical comparison, and explicit claim-status analysis.
%
% Foundational Restrictions :
%   Heidegger's distinction between Being and beings is not
%   equated with the distinction between Being and Identity.
%   Dasein's understanding of Being is not made the ontological
%   source of Being or Identity.
%   Heidegger's existential structures are not imported as
%   foundational definitions of the ontology developed in Book I.
%   Lichtung is not identified with Ontological Light.
%
% Claim Form and Epistemic Status :
%   Historical reconstructions are textually grounded.
%   Interpretive disputes are explicitly qualified.
%   Comparative conclusions are argued assessments.
%   Unresolved questions remain open and are not promoted to
%   foundational commitments.
%
% =============================================================================

\section{Martin Heidegger}
\label{sec:ch08_s11}

\index[name]{Heidegger, Martin}
\index[concept]{Ontological Difference}
\index[concept]{Dasein}
\index[concept]{Being-in-the-World}
\index[concept]{Worldhood}
\index[concept]{Readiness-to-Hand}
\index[concept]{Presence-at-Hand}
\index[concept]{Being-with}
\index[concept]{das Man}
\index[concept]{Care}
\index[concept]{Being-toward-Death}
\index[concept]{Authenticity}
\index[concept]{Temporality}
\index[concept]{Disclosedness}
\index[concept]{Lichtung}
\index[concept]{Kehre}
\index[concept]{Ereignis}
\index[concept]{Enframing}
\index[concept]{Dwelling}


% =============================================================================
% 1. SCOPE AND COMPARATIVE ORIENTATION
% =============================================================================

\subsection{Scope and Comparative Orientation}
\label{subsec:ch08-heidegger-scope}

% Structural ID: VS.B01.P01.C08.S11.SS01.P001
\phantomsection\label{para:VS.B01.P01.C08.S11.SS01.P001}

Martin Heidegger is a decisive interlocutor for an ontology
concerned with Being and Identity because he challenges the
assumption that philosophical inquiry can begin and end with
the classification of entities, their properties, or their
relations. His question concerns Being itself: what it means
for entities to be, and how their Being becomes intelligible.
The force of this question is not exhausted by the distinction
between one entity and another. It reaches the ontological
horizon within which entities can be encountered as entities.

% Structural ID: VS.B01.P01.C08.S11.SS01.P002
\phantomsection\label{para:VS.B01.P01.C08.S11.SS01.P002}

The present comparison concentrates on the relation between
Heidegger's fundamental ontology and the distinctions
developed in Part I. Its principal subjects are the
ontological difference, Dasein, being-in-the-world, worldhood,
equipmental involvement, being-with, care, authenticity,
temporality, truth as unconcealment, and the later thinking
of Being. These subjects are examined because they place
pressure on any ontology that speaks of identity,
determination, manifestation, and relational constitution
without first clarifying the meaning of Being.

% Structural ID: VS.B01.P01.C08.S11.SS01.P003
\phantomsection\label{para:VS.B01.P01.C08.S11.SS01.P003}

The comparison is not an attempt to translate Heidegger's
vocabulary into the terminology of the Philosophy of Viability.
Being is not a synonym for Identity; Dasein is not a synonym
for an identity-bearing entity; worldhood is not simply
Structure; and unconcealment is not simply Manifestation.
The comparison must establish where a genuine conceptual
relation exists and where a resemblance conceals a difference
in philosophical function. The methodological requirements
established in the opening section of this chapter govern
the present examination.

% Structural ID: VS.B01.P01.C08.S11.SS01.P004
\phantomsection\label{para:VS.B01.P01.C08.S11.SS01.P004}

The central thesis is that Heidegger reveals a dimension
of ontological intelligibility that a theory of identity
cannot simply bypass: entities are encountered within
horizons of significance, practical involvement, social
relations, and temporality. This is a substantial challenge
to any account that treats entities as self-contained
objects whose relations are merely external additions.
Yet Heidegger's analytic of Dasein does not, by itself,
supply a complete ontology of Identity. The relation
between the question of Being and the ontological status
of Identity therefore remains a problem for critical
comparison rather than a conclusion to be assumed.


% =============================================================================
% 2. HISTORICAL AND TEXTUAL LOCATION
% =============================================================================

\subsection{Historical and Textual Location}
\label{subsec:ch08-heidegger-historical}

% Structural ID: VS.B01.P01.C08.S11.SS02.P001
\phantomsection\label{para:VS.B01.P01.C08.S11.SS02.P001}

Heidegger's principal early work, \textit{Being and Time},
was published in 1927. Its declared project is to reopen
the question of the meaning of Being and to prepare that
inquiry through an existential analytic of Dasein. The
published work does not contain the complete projected
treatise: the intended further development of the question
of Being was not published in the form originally announced.
This incompleteness matters because the analytic of Dasein
must not be mistaken for the completed exposition of every
dimension of Heidegger's ontology.
\cite[Introduction and §§1--8]{HeideggerBeingTime}

% Structural ID: VS.B01.P01.C08.S11.SS02.P002
\phantomsection\label{para:VS.B01.P01.C08.S11.SS02.P002}

The early analysis proceeds through the structures of
existence, being-in-the-world, worldhood, being-with,
everydayness, care, disclosedness, death, conscience,
resoluteness, and temporality. Heidegger's later writings
develop the question through a different vocabulary,
including the history of Being, unconcealment, the
clearing, Ereignis, technology, language, and dwelling.
The early and later works are connected by recurring
questions, but their relation cannot be settled by
treating the later vocabulary as a simple continuation
of the terminology of 1927.

% Structural ID: VS.B01.P01.C08.S11.SS02.P003
\phantomsection\label{para:VS.B01.P01.C08.S11.SS02.P003}

The interpretation of the relation between early and
later Heidegger remains disputed. Some readings emphasize
continuity in the question of Being and the interplay
between disclosure and concealment; others emphasize
changes in the account of ontological dependence, the
historicity of Being, and the role assigned to human
existence. The term \textit{Kehre} is useful for identifying
this development, but it should not be treated as proof
that one sharply defined doctrinal reversal occurred
at a single moment.
\footnote{For a recent reconstruction of the early and
later phases, including the dispute over the nature and
extent of the \textit{Kehre}, see Mark Wrathall,
``Martin Heidegger,'' \textit{The Stanford Encyclopedia
of Philosophy}, Spring 2025 Edition. \cite{WrathallHeideggerSEP2025}}

% Structural ID: VS.B01.P01.C08.S11.SS02.P004
\phantomsection\label{para:VS.B01.P01.C08.S11.SS02.P004}

The present section gives priority to \textit{Being and Time}
for the existential analytic and uses selected later
writings to examine the development of disclosure,
technology, and dwelling. It does not attempt a complete
history of Heidegger's writings. This restriction is
important because a claim about the early analytic cannot
automatically be attributed to the later thought of Being,
and a later formulation cannot be retroactively inserted
into the argument of the 1927 text without qualification.


% =============================================================================
% 3. THE QUESTION OF BEING AND THE ONTOLOGICAL DIFFERENCE
% =============================================================================

\subsection{The Question of Being and the Ontological Difference}
\label{subsec:ch08-heidegger-ontological-difference}

% Structural ID: VS.B01.P01.C08.S11.SS03.P001
\phantomsection\label{para:VS.B01.P01.C08.S11.SS03.P001}

Heidegger begins from a distinction between asking what
an entity is and asking what it means for an entity to be.
The first question may identify a kind, a property, a cause,
or a constitutive feature. The second asks about the
intelligibility of Being through which entities can be
encountered and understood as entities. The question of
Being is therefore not another question about a particularly
important entity. It concerns the ontological presupposition
that is already operative whenever entities are investigated.
\cite[Introduction, §§1--4]{HeideggerBeingTime}

% Structural ID: VS.B01.P01.C08.S11.SS03.P002
\phantomsection\label{para:VS.B01.P01.C08.S11.SS03.P002}

The ontological difference names the distinction between
Being (\textit{Sein}) and beings (\textit{Seiendes}).
Being is not one more being alongside other beings, nor
is it a highest entity whose properties could be listed
within the same classificatory order. Conversely, beings
are not reducible to the concept of Being. The distinction
concerns two different levels of inquiry: the entities
that are, and the Being in terms of which they are
understood as entities. Heidegger's later writings
make this difference increasingly explicit.

% Structural ID: VS.B01.P01.C08.S11.SS03.P003
\phantomsection\label{para:VS.B01.P01.C08.S11.SS03.P003}

The distinction has a critical consequence. A philosophical
account may provide an extensive classification of beings
while leaving unexamined the meaning of Being presupposed
by that classification. The accumulation of categories
does not automatically answer the question of Being.
Heidegger's criticism is therefore directed not merely
against particular answers but against the possibility
that ontology may mistake the articulation of entities
for the clarification of Being itself.
\cite[§§2--4]{HeideggerBeingTime}

% Structural ID: VS.B01.P01.C08.S11.SS03.P004
\phantomsection\label{para:VS.B01.P01.C08.S11.SS03.P004}

The comparison with Book I must preserve the difference
between Heidegger's ontological difference and the
analytical distinction between Being and Identity.
The former distinguishes Being from beings; the latter
distinguishes two ontological dimensions that this book
holds to be inseparable and without ontological priority
over one another. These distinctions address different
questions. Heidegger's distinction cannot be substituted
for the three independent principles of Analytical
Distinction, Ontological Inseparability, and Absence
of Ontological Priority.

% Structural ID: VS.B01.P01.C08.S11.SS03.P005
\phantomsection\label{para:VS.B01.P01.C08.S11.SS03.P005}

The strongest challenge nevertheless remains. If an
ontology develops an account of Identity, Determination,
and Manifestation without addressing the meaning of Being,
does it risk describing beings while leaving their Being
unexamined? This question cannot be dismissed by repeating
that Book I is an ontology. It requires an argument about
the scope of the book's inquiry and about whether its
foundational distinctions presuppose a conception of
Being that has not yet been sufficiently clarified.
\footnote{The scope and incompleteness of the project
of \textit{Being and Time}, and the interpretive problem
of its fundamental-ontological ambition, are discussed
in Wrathall, ``Martin Heidegger.'' \cite{WrathallHeideggerSEP2025}}


% =============================================================================
% 4. DASEIN AND THE EXISTENTIAL ANALYTIC
% =============================================================================

\subsection{Dasein and the Existential Analytic}
\label{subsec:ch08-heidegger-dasein}

% Structural ID: VS.B01.P01.C08.S11.SS04.P001
\phantomsection\label{para:VS.B01.P01.C08.S11.SS04.P001}

Heidegger introduces Dasein as the entity that we ourselves
are and whose Being is an issue for it. Dasein is not
introduced as a psychological subject, a biological
organism considered only as a living body, or a Cartesian
consciousness placed before an external world. The term
marks an ontological inquiry into the kind of Being
belonging to human existence. Dasein is distinguished
by an understanding of Being that is operative in its
existence, even when that understanding has not become
an explicit philosophical theory.
\cite[§§4, 9--10]{HeideggerBeingTime}

% Structural ID: VS.B01.P01.C08.S11.SS04.P002
\phantomsection\label{para:VS.B01.P01.C08.S11.SS04.P002}

The existential analytic investigates the structures
that make Dasein's existence intelligible. Heidegger
distinguishes existential structures from particular
existentiell ways in which those structures are lived.
The distinction is methodological and ontological:
an existential describes a structural feature of
existence, while an existentiell determination concerns
a concrete way in which an individual lives or takes
up a possibility. The analytic is not an empirical
inventory of psychological traits.

% Structural ID: VS.B01.P01.C08.S11.SS04.P003
\phantomsection\label{para:VS.B01.P01.C08.S11.SS04.P003}

Dasein is essentially related to possibilities. It
does not first possess a complete and fixed definition
and then happen to choose among optional activities.
Its existence involves projection into possibilities
that it can take up, avoid, or fail to recognize.
This modal structure explains why Heidegger resists
defining Dasein solely through a list of present-at-hand
properties. The question of who Dasein is cannot be
separated from the possibilities through which it
understands and conducts its existence.
\cite[§§9, 31, 41]{HeideggerBeingTime}

% Structural ID: VS.B01.P01.C08.S11.SS04.P004
\phantomsection\label{para:VS.B01.P01.C08.S11.SS04.P004}

This analysis is significant for an ontology of identity
because it exposes the inadequacy of treating a human
being as a closed collection of currently instantiated
properties. Yet the existential structure of Dasein
must not be generalized into a universal definition
of every existent. Nor does the fact that Dasein
understands Being establish that Being or Identity
depends ontologically upon human understanding.
The ontological status of non-Dasein entities remains
a question that cannot be settled merely by extending
the existential analytic beyond its stated domain.


% =============================================================================
% 5. BEING-IN-THE-WORLD AND WORLDHOOD
% =============================================================================

\subsection{Being-in-the-World and Worldhood}
\label{subsec:ch08-heidegger-worldhood}

% Structural ID: VS.B01.P01.C08.S11.SS05.P001
\phantomsection\label{para:VS.B01.P01.C08.S11.SS05.P001}

The fundamental structure of Dasein is being-in-the-world
(\textit{In-der-Welt-sein}). The hyphenated expression
does not describe the spatial inclusion of a subject
inside a container called the world. It names a unitary
phenomenon in which Dasein, world, and being-in belong
together. The world is not first encountered as a
neutral collection of objects to which a subject later
assigns meanings. It is already encountered through
practical involvements, shared activities, and
referential relations.
\cite[§§12--14]{HeideggerBeingTime}

% =============================================================================
% Project : THE PHILOSOPHY OF VIABILITY
% Book    : Book I — Ontology of Viability
% Book ID : VS.B01
% Part    : Part I — Identity and Manifestation
% Part ID : VS.B01.P01
% Chapter : Chapter 08 — Comparative Philosophical Examination
% Chapter ID : VS.B01.P01.C08
%
% File :
%   4_ASSETS/illustrations/
%   VS.B01.P01.C08.S11.SS05.P001.IL01.heidegger_equipmental_world.tex
%
% Section : Section 11 — Martin Heidegger
%
% Asset :
%   Illustration IL01 — Heidegger's Equipmental World
%
% Structural Identifier :
%   VS.B01.P01.C08.S11.SS05.P001.IL01
%
% Purpose :
%   To provide an illustrative scene for Heidegger's analysis
%   of equipment, practical involvement, and worldhood.
%
% Intellectual Benefit :
%   Makes the referential context of equipment approachable.
%
% Aesthetic / Experiential Benefit :
%   Provides a visual pause and a memorable scene within
%   a conceptually dense discussion.
%
% Representational Limit :
%   This is an illustration of a practical situation.
%   It is not a complete diagram of Heidegger's ontology,
%   nor does it assert that all entities derive their Being
%   from human use.
%
% Dependencies :
%   TikZ; colors defined in the project preamble.
%
% =============================================================================

\begin{figure}[htbp]
\centering

\begin{tikzpicture}[x=1cm,y=1cm]

    % Background
    \fill[ViabilityPaper] (0,0) rectangle (13,6.2);

    % Workshop wall and floor
    \fill[ViabilityDarkBlue!5] (0,0) rectangle (13,1.15);
    \draw[ViabilityDarkBlue!35,line width=0.8pt]
        (0,1.15) -- (13,1.15);

    % Workbench
    \fill[ViabilityDark!18]
        (1.0,1.25) rectangle (11.9,1.65);
    \draw[ViabilityDark!70,line width=1pt]
        (1.0,1.25) rectangle (11.9,1.65);
    \draw[ViabilityDark!70,line width=1.2pt]
        (1.7,1.25) -- (1.7,0.25)
        (11.2,1.25) -- (11.2,0.25);

    % Timber / board
    \fill[ViabilityGold!28]
        (3.3,1.72) rectangle (9.4,2.12);
    \draw[ViabilityGold!75!black,line width=1pt]
        (3.3,1.72) rectangle (9.4,2.12);

    % Nail
    \draw[ViabilityDark,line width=1.4pt]
        (6.45,2.12) -- (6.45,2.65);
    \draw[ViabilityDark,line width=2pt]
        (6.30,2.65) -- (6.60,2.65);

    % Hammer head
    \fill[ViabilityGray!35]
        (5.35,3.55) rectangle (7.2,4.05);
    \draw[ViabilityDark,line width=1.1pt]
        (5.35,3.55) rectangle (7.2,4.05);
    \draw[ViabilityDark,line width=1.1pt]
        (5.35,3.55) -- (5.05,3.55)
        (5.05,3.55) -- (5.05,4.05)
        (5.05,4.05) -- (5.35,4.05);

    % Hammer handle
    \fill[ViabilityGold!55]
        (6.0,3.55) -- (6.45,3.55)
        -- (6.9,2.55) -- (6.45,2.40) -- cycle;
    \draw[ViabilityGold!80!black,line width=1pt]
        (6.0,3.55) -- (6.45,3.55)
        -- (6.9,2.55) -- (6.45,2.40) -- cycle;

    % Workshop window
    \fill[ViabilityBlue!8]
        (9.6,3.2) rectangle (12.1,5.3);
    \draw[ViabilityDarkBlue!55,line width=1pt]
        (9.6,3.2) rectangle (12.1,5.3);
    \draw[ViabilityDarkBlue!40,line width=0.7pt]
        (10.85,3.2) -- (10.85,5.3)
        (9.6,4.25) -- (12.1,4.25);

    % Outline of a house under construction
    \draw[ViabilityDarkBlue!65,line width=1.2pt]
        (0.9,3.0) -- (2.1,4.1) -- (3.3,3.0)
        (1.25,3.0) -- (1.25,2.15)
        (2.95,3.0) -- (2.95,2.15)
        (1.25,2.15) -- (2.95,2.15);

    % Referential connections
    \draw[ViabilityBlue!75!black,dashed,line width=0.9pt]
        (6.35,4.15) .. controls (5.2,5.0) and (3.8,4.8) .. (2.2,4.15);
    \draw[ViabilityBlue!75!black,dashed,line width=0.9pt]
        (7.15,3.8) .. controls (8.0,4.8) and (9.0,4.8) .. (10.2,4.7);
    \draw[ViabilityBlue!75!black,dashed,line width=0.9pt]
        (6.45,2.45) -- (6.45,2.12);

    % Labels
    \node[
        anchor=west,
        text=ViabilityDarkBlue,
        font=\large\bfseries
    ] at (0.55,5.75) {A Workshop and Its World};

    \node[
        anchor=west,
        text=ViabilityGray,
        font=\small
    ] at (0.58,5.35)
        {An illustrative scene of practical involvement};

    \node[
        text=ViabilityDarkBlue,
        font=\small\bfseries
    ] at (2.1,1.95) {Building};

    \node[
        text=ViabilityDarkBlue,
        font=\small\bfseries
    ] at (6.4,4.45) {Tool};

    \node[
        text=ViabilityDarkBlue,
        font=\small\bfseries
    ] at (10.9,2.85) {Workshop};

    \node[
        text=ViabilityGray,
        font=\footnotesize
    ] at (6.5,0.55)
        {Task \quad Material \quad Equipment \quad Place};

\end{tikzpicture}

\caption{An illustrative workshop scene inspired by Heidegger's
analysis of equipment and worldhood. The hammer is presented
within a practical context involving materials, tasks, tools,
and a place of work. The illustration is not a complete
representation of Heidegger's existential analytic.}

\label{fig:VS.B01.P01.C08.S11.SS05.P001.IL01}

\end{figure}

% =============================================================================
% END OF ILLUSTRATION IL01
% =============================================================================

% Structural ID: VS.B01.P01.C08.S11.SS05.P002
\phantomsection\label{para:VS.B01.P01.C08.S11.SS05.P002}

Worldhood is disclosed through networks of involvement.
A hammer refers to nails, materials, a task, a workshop,
and a project; these relations are not merely accidental
associations added to an otherwise complete object.
They belong to the context in which the equipment can
be encountered as useful equipment. The world is
therefore not an aggregate of entities. It is the
referential totality within which particular entities
become intelligible in their practical roles.
\cite[§§15--18]{HeideggerBeingTime}

% Structural ID: VS.B01.P01.C08.S11.SS05.P003
\phantomsection\label{para:VS.B01.P01.C08.S11.SS05.P003}

This analysis challenges an ontology that treats relations
as merely external connections between independently
complete objects. Heidegger shows that at least some
modes of encounter depend upon a context of involvement
that is not reducible to the sum of its components.
The point is not that every entity is constituted by
human use, but that the Being of equipment cannot be
adequately described by abstracting it from the
referential context in which it is encountered.

% Structural ID: VS.B01.P01.C08.S11.SS05.P004
\phantomsection\label{para:VS.B01.P01.C08.S11.SS05.P004}

For Book I, the analysis opens a question about the
relation between relational determination and the
ontological status of an existent. It gives reason
to examine whether certain determinations are
intelligible only within relational contexts. It does
not establish that Identity is identical with those
relations, that every being is constituted by one
world of significance, or that the ontology of all
beings can be reduced to the practical world of Dasein.
\footnote{Wrathall's reconstruction of worldhood
emphasizes the holistic and practical structure of
being-in-the-world while also identifying questions
about the material body and spatiality. \cite{WrathallHeideggerSEP2025}}


% =============================================================================
% 6. EQUIPMENT AND MODES OF ENCOUNTER
% =============================================================================

\subsection{Equipment and Modes of Encounter}
\label{subsec:ch08-heidegger-equipment}

% Structural ID: VS.B01.P01.C08.S11.SS06.P001
\phantomsection\label{para:VS.B01.P01.C08.S11.SS06.P001}

Heidegger distinguishes readiness-to-hand
(\textit{Zuhandenheit}) from presence-at-hand
(\textit{Vorhandenheit}). Equipment is ordinarily
encountered in the course of practical activity:
the hammer is used in hammering rather than first
being inspected as an object with a set of properties.
In this mode, the tool is encountered through its
involvement in a task and through the network of
references that makes the task possible.
\cite[§§15--16]{HeideggerBeingTime}

% Structural ID: VS.B01.P01.C08.S11.SS06.P002
\phantomsection\label{para:VS.B01.P01.C08.S11.SS06.P002}

When equipment breaks, is missing, or obstructs an
activity, its ordinary practical transparency may
be interrupted. The tool can then become conspicuous
as an object of inspection. This change in the mode
of encounter reveals the difference between practical
involvement and detached observation. Presence-at-hand
is not simply false, and readiness-to-hand is not
a subjective illusion; they are distinct ways in
which entities can be encountered.
\cite[§§16--18]{HeideggerBeingTime}

% Structural ID: VS.B01.P01.C08.S11.SS06.P003
\phantomsection\label{para:VS.B01.P01.C08.S11.SS06.P003}

The distinction exposes a limitation in the assumption
that the most fundamental description of an entity
must be an inventory of properties available to
detached observation. Practical involvement can
disclose relations that are obscured when an entity
is considered only as an isolated object. The
analysis therefore gives philosophical significance
to modes of encounter without reducing the entity
to the act of encountering it.

% Structural ID: VS.B01.P01.C08.S11.SS06.P004
\phantomsection\label{para:VS.B01.P01.C08.S11.SS06.P004}

The comparison must nevertheless distinguish the
mode in which an entity is encountered from the
identity of that entity. A change from readiness-to-hand
to presence-at-hand does not by itself establish
a change of Identity. Nor does the interruption
of practical involvement prove that the entity has
lost its Being. Heidegger's analysis concerns modes
of Being and disclosure; the ontological conditions
of identity through such changes require a further
argument.


% =============================================================================
% 7. BEING-WITH AND EVERYDAYNESS
% =============================================================================

\subsection{Being-with and Everydayness}
\label{subsec:ch08-heidegger-being-with}

% Structural ID: VS.B01.P01.C08.S11.SS07.P001
\phantomsection\label{para:VS.B01.P01.C08.S11.SS07.P001}

Dasein is essentially being-with (\textit{Mitsein}).
Others are not first encountered as isolated objects
from which a subject subsequently infers the existence
of other minds. Shared practices, equipment, norms,
and activities already disclose a world in which
others are involved. Being-with is therefore an
existential structure of Dasein, not merely an
empirical fact that several individuals happen
to occupy the same location.
\cite[§§25--27]{HeideggerBeingTime}

% Structural ID: VS.B01.P01.C08.S11.SS07.P002
\phantomsection\label{para:VS.B01.P01.C08.S11.SS07.P002}

The analysis of everydayness introduces das Man,
often translated as ``the they'' or ``the anyone''.
The expression identifies the anonymous public
norms through which everyday possibilities are
interpreted and distributed. It is not simply
a crowd of other people. It names the impersonal
``one'' that determines what one ordinarily does,
says, expects, and considers appropriate.
\cite[§27]{HeideggerBeingTime}

% Structural ID: VS.B01.P01.C08.S11.SS07.P003
\phantomsection\label{para:VS.B01.P01.C08.S11.SS07.P003}

The they has a double significance. It can level
possibilities and encourage conformity, but it
also sustains the shared intelligibility required
for ordinary social existence. Heidegger's account
therefore cannot be reduced to a moral denunciation
of society or to the claim that individuality is
always destroyed by common practices. Everyday
sociality is a constitutive condition of the world
in which Dasein understands itself.
\cite[§§27, 35--38]{HeideggerBeingTime}

% Structural ID: VS.B01.P01.C08.S11.SS07.P004
\phantomsection\label{para:VS.B01.P01.C08.S11.SS07.P004}

For the ontology of identity, the analysis raises
a question about the relation between individual
identity and socially shared determinations. A
person's social roles and practical relations may
be significant to the way that person exists and
is understood, but the account of those relations
does not by itself determine the number or nature
of the person's identities. The plurality of roles
must not be confused with a plurality of identities.
The latter remains a separate ontological problem.


% =============================================================================
% 8. CARE, THROWNNESS, AND PROJECTION
% =============================================================================

\subsection{Care, Thrownness, and Projection}
\label{subsec:ch08-heidegger-care}

% Structural ID: VS.B01.P01.C08.S11.SS08.P001
\phantomsection\label{para:VS.B01.P01.C08.S11.SS08.P001}

Heidegger's analysis of care (\textit{Sorge}) seeks to
grasp Dasein as a structural whole. Care is not
primarily an emotion, a moral virtue, or a voluntary
attitude toward another person. It names the unity
of Dasein's being-ahead-of-itself, its already-being-in
a world, and its being-alongside entities encountered
within that world. These moments belong together
and cannot be treated as independent components
that are later assembled into existence.
\cite[§§39--42]{HeideggerBeingTime}

% Structural ID: VS.B01.P01.C08.S11.SS08.P002
\phantomsection\label{para:VS.B01.P01.C08.S11.SS08.P002}

Being-ahead-of-itself concerns projection into
possibilities. Already-being-in expresses thrownness:
Dasein always finds itself within conditions and
histories that it did not first choose. Being-alongside
concerns its involvement with entities and tasks
within the world. Care holds these dimensions together
as the formal structure of Dasein's existence.
It is not a causal sequence in which one moment
produces the next.
\cite[§41]{HeideggerBeingTime}

% Structural ID: VS.B01.P01.C08.S11.SS08.P003
\phantomsection\label{para:VS.B01.P01.C08.S11.SS08.P003}

The structure of care also clarifies why Dasein cannot
be understood exclusively through either facticity
or possibility. It is thrown into a world, yet it
projects possibilities from within that thrown
situation. It is involved with entities, yet its
own Being is not exhausted by the entities with
which it deals. Heidegger's analysis therefore
connects situatedness, projection, and practical
involvement without reducing one to another.

% Structural ID: VS.B01.P01.C08.S11.SS08.P004
\phantomsection\label{para:VS.B01.P01.C08.S11.SS08.P004}

The following diagram presents the internal relations
of the existential analytic. It is an interpretive
map of Heidegger's account, not a causal model and
not a general architecture of all beings. In particular,
the unity of care and the temporal interpretation
of that unity must not be mistaken for a definition
of Identity or for an established structure of
the Ontology of Viability.

% =============================================================================
% Project : THE PHILOSOPHY OF VIABILITY
% Book    : Book I — Ontology of Viability
% Book ID : VS.B01
% Part    : Part I — Identity and Manifestation
% Part ID : VS.B01.P01
% Chapter : Chapter 08 — Comparative Philosophical Examination
% Chapter ID : VS.B01.P01.C08
%
% File :
%   4_ASSETS/diagrams/
%   VS.B01.P01.C08.S11.SS08.P004.DG01.heidegger_dasein_care_and_temporality.tex
%
% Section : Section 11 — Martin Heidegger
%
% Asset :
%   Diagram DG01 — Dasein, Care, and Temporality
%
% Structural Identifier :
%   VS.B01.P01.C08.S11.SS08.P004.DG01
%
% Purpose :
%   To represent the structural relations among
%   being-in-the-world, care, and ecstatic temporality.
%
% Intellectual Benefit :
%   Clarifies the internal organization of Heidegger's
%   existential analytic and prevents a linear reading
%   of its interdependent structures.
%
% Aesthetic / Experiential Benefit :
%   Provides a compact visual synthesis of a dense argument.
%
% Representational Limit :
%   The diagram is an interpretive map of Dasein.
%   It is not a causal sequence or an ontology of all beings.
%
% Dependencies :
%   TikZ; positioning library; project colors.
%
% =============================================================================

\begin{figure}[htbp]
\centering

\begin{tikzpicture}[
    font=\small,
    node distance=0.55cm,
    structure/.style={
        draw=ViabilityBlue!75!black,
        fill=ViabilityPaleBlue,
        rounded corners=3pt,
        line width=0.7pt,
        text width=3.5cm,
        minimum height=1.0cm,
        align=center,
        inner sep=5pt
    },
    central/.style={
        draw=ViabilityGold!85!black,
        fill=ViabilityPaleGold,
        rounded corners=3pt,
        line width=1pt,
        text width=10.6cm,
        minimum height=0.9cm,
        align=center,
        font=\small\bfseries,
        inner sep=6pt
    },
    temporal/.style={
        draw=ViabilityGreen!75!black,
        fill=ViabilityPaleGreen,
        rounded corners=3pt,
        line width=0.8pt,
        text width=3.1cm,
        minimum height=0.85cm,
        align=center,
        inner sep=5pt
    },
    limit/.style={
        draw=ViabilityGray,
        fill=ViabilityLightGray!35,
        rounded corners=3pt,
        text width=10.6cm,
        align=center,
        font=\footnotesize,
        inner sep=6pt
    },
    relation/.style={
        draw=ViabilityDarkBlue,
        line width=0.8pt,
        -{Stealth[length=2mm]}
    }
]

\node[structure] (world) {
    \textbf{Worldhood}\\
    Referential context and significance
};

\node[structure,right=0.35cm of world] (who) {
    \textbf{Who}\\
    Being-with and the anyone-self
};

\node[structure,right=0.35cm of who] (beingin) {
    \textbf{Being-in}\\
    Disposedness, understanding, discourse
};

\node[central,below=0.8cm of who] (whole) {
    Being-in-the-world\\
    A unitary existential phenomenon
};

\draw[relation] (world.south) -- (whole.north west);
\draw[relation] (who.south) -- (whole.north);
\draw[relation] (beingin.south) -- (whole.north east);

\node[central,below=0.65cm of whole] (care) {
    Care (\textit{Sorge}) — the structural whole of Dasein
};

\draw[relation] (whole) -- (care);

\node[structure,below=0.55cm of care,xshift=-3.55cm] (projection) {
    \textbf{Being-ahead-of-itself}\\
    Projection into possibilities
};

\node[structure,below=0.55cm of care] (thrown) {
    \textbf{Already-being-in}\\
    Thrownness and facticity
};

\node[structure,below=0.55cm of care,xshift=3.55cm] (involvement) {
    \textbf{Being-alongside}\\
    Involvement with entities
};

\draw[relation] (care.south) -- (projection.north);
\draw[relation] (care.south) -- (thrown.north);
\draw[relation] (care.south) -- (involvement.north);

\node[central,below=0.8cm of thrown] (temporality) {
    Temporality — the ontological meaning of care
};

\draw[relation] (projection.south) -- (temporality.north west);
\draw[relation] (thrown.south) -- (temporality.north);
\draw[relation] (involvement.south) -- (temporality.north east);

\node[temporal,below=0.55cm of temporality,xshift=-3.55cm] (future) {
    \textbf{Future}\\
    Coming-toward
};

\node[temporal,below=0.55cm of temporality] (past) {
    \textbf{Having-been}\\
    Inherited situation
};

\node[temporal,below=0.55cm of temporality,xshift=3.55cm] (present) {
    \textbf{Present}\\
    Making-present
};

\draw[relation] (temporality.south) -- (future.north);
\draw[relation] (temporality.south) -- (past.north);
\draw[relation] (temporality.south) -- (present.north);

\node[limit,below=0.75cm of past] (boundary) {
    \textbf{Comparative limit}\\
    This structure concerns Dasein's existential analytic.
    It is not a causal model, a universal inventory of beings,
    or a definition of Identity in the Ontology of Viability.
};

\end{tikzpicture}

\caption{The existential structure of Dasein in Heidegger's
analysis. Being-in-the-world is interpreted through care,
whose temporal meaning is articulated through the ecstatic
unity of future, having-been, and present. The relations
are structural rather than a sequence of efficient causes.}

\label{fig:VS.B01.P01.C08.S11.SS08.P004.DG01}

\end{figure}

% =============================================================================
% END OF DIAGRAM DG01
% =============================================================================


% =============================================================================
% 9. BEING-TOWARD-DEATH, AUTHENTICITY, AND SELFHOOD
% =============================================================================

\subsection{Being-toward-Death, Authenticity, and Selfhood}
\label{subsec:ch08-heidegger-authenticity}

% Structural ID: VS.B01.P01.C08.S11.SS09.P001
\phantomsection\label{para:VS.B01.P01.C08.S11.SS09.P001}

Being-toward-death (\textit{Sein-zum-Tode}) is not simply
the anticipation of a future biological event. It
concerns Dasein's relation to its ownmost possibility
of no-longer-being-able-to-exist. Death is certain
but indeterminate in its when, and no one can take
over another's dying in the existential sense
developed by Heidegger. Anticipation of death
individualizes Dasein by disclosing the finitude
of its possibilities.
\cite[§§46--53]{HeideggerBeingTime}

% Structural ID: VS.B01.P01.C08.S11.SS09.P002
\phantomsection\label{para:VS.B01.P01.C08.S11.SS09.P002}

Authenticity (\textit{Eigentlichkeit}) concerns the
way Dasein takes over its own possibilities rather
than simply allowing the anonymous norms of everyday
life to determine them. Inauthenticity is not an
ontological defect that could be eliminated once
and for all. It is an ordinary mode of existence
that belongs to Dasein's everyday involvement.
Authenticity does not disclose a hidden substance
behind social existence; it concerns how existence
is owned and enacted.
\cite[§§54--60]{HeideggerBeingTime}

% Structural ID: VS.B01.P01.C08.S11.SS09.P003
\phantomsection\label{para:VS.B01.P01.C08.S11.SS09.P003}

The call of conscience and resoluteness belong to
this analysis. The call does not provide a catalogue
of moral rules or additional information about
worldly events. It summons Dasein to take responsibility
for its own finite possibilities. Resoluteness is
therefore not a guarantee of moral correctness;
it is an existential mode in which Dasein takes
over its thrown possibilities.
\cite[§§54--60]{HeideggerBeingTime}

% Structural ID: VS.B01.P01.C08.S11.SS09.P004
\phantomsection\label{para:VS.B01.P01.C08.S11.SS09.P004}

The analysis of selfhood is important but limited
for the present ontology. Heidegger examines the
constancy of the self through anticipatory resoluteness,
but this is not equivalent to a complete account
of numerical identity, ontological individuation,
or persistence through every kind of transformation.
The existential meaning of owning one's possibilities
must remain distinct from the ontological question
of what makes an existent the existent that it is.
\footnote{Interpretations of authenticity differ over
the role of social norms, individualization, and
responsibility. Wrathall surveys these interpretive
issues in the sections on authenticity and death.
\cite{WrathallHeideggerSEP2025}}


% =============================================================================
% 10. TEMPORALITY AND HISTORICALITY
% =============================================================================

\subsection{Temporality and Historicality}
\label{subsec:ch08-heidegger-temporality}

% Structural ID: VS.B01.P01.C08.S11.SS10.P001
\phantomsection\label{para:VS.B01.P01.C08.S11.SS10.P001}

Heidegger's account of temporality
(\textit{Zeitlichkeit}) is not a theory of time as
a neutral container in which entities move from
one present moment to another. Temporality is the
ecstatic unity of future, having-been, and present.
These are not three successive pieces of time.
They are interdependent dimensions through which
Dasein projects possibilities, inherits a thrown
situation, and makes present the entities with
which it is involved.
\cite[§§65--71]{HeideggerBeingTime}

% Structural ID: VS.B01.P01.C08.S11.SS10.P002
\phantomsection\label{para:VS.B01.P01.C08.S11.SS10.P002}

The future has a distinctive role because Dasein
exists through possibilities toward which it
projects itself. Yet this projection is never
detached from having-been: possibilities are
taken up from a situation and a history that
Dasein did not create from nothing. The present
is correspondingly understood through the
activities and concerns that become intelligible
within this temporal unity. Temporality is thus
not a linear succession but a structured unity
of ecstatic dimensions.

% Structural ID: VS.B01.P01.C08.S11.SS10.P003
\phantomsection\label{para:VS.B01.P01.C08.S11.SS10.P003}

Heidegger's project also connects temporality
with the meaning of Being. The analysis of
Dasein is intended to disclose the temporal
horizon within which Being can be understood.
The later concern with the history of Being
develops this issue by examining historical
transformations in the prevailing ways entities
are disclosed. The relation between existential
temporality and the history of Being remains
an important point of development and dispute.
\footnote{Wrathall distinguishes Dasein's
\textit{Zeitlichkeit}, world-time, and the broader
question of \textit{Temporalität}, while also
discussing the historical dimension of the later
thought. \cite{WrathallHeideggerSEP2025}}

% Structural ID: VS.B01.P01.C08.S11.SS10.P004
\phantomsection\label{para:VS.B01.P01.C08.S11.SS10.P004}

For Book I, temporality provides a serious challenge
to any account of continuity that treats time as
an external sequence of states. It directs attention
to the temporal conditions under which possibilities,
relations, and modes of manifestation become
intelligible. It does not establish that Identity
is identical with temporal continuity, that every
identity is constituted by Dasein's temporality,
or that the temporal unity of existence proves
numerical identity through all transformations.


% =============================================================================
% 11. TRUTH, DISCLOSURE, AND THE CLEARING
% =============================================================================

\subsection{Truth, Disclosure, and the Clearing}
\label{subsec:ch08-heidegger-clearing}

% Structural ID: VS.B01.P01.C08.S11.SS11.P001
\phantomsection\label{para:VS.B01.P01.C08.S11.SS11.P001}

Heidegger re-examines truth through the Greek
notion of \textit{aletheia}, commonly rendered
as unconcealment or unhiddenness. The claim is
not that propositional correctness becomes
meaningless. Rather, correctness is treated as
a derivative possibility that presupposes a
more primordial openness in which entities
can first be disclosed as the entities about
which propositions may be true or false.
\cite[§44]{HeideggerBeingTime}

% Structural ID: VS.B01.P01.C08.S11.SS11.P002
\phantomsection\label{para:VS.B01.P01.C08.S11.SS11.P002}

The later notion of the clearing
(\textit{Lichtung}) develops the image of an
open region in which entities can appear and
remain concealed. Disclosure and concealment
belong together: the opening of one field of
intelligibility does not mean that everything
is revealed exhaustively. The clearing is
therefore not simply a luminous object, a
physical light, or a transparent space in
which all entities become completely known.

% Structural ID: VS.B01.P01.C08.S11.SS11.P003
\phantomsection\label{para:VS.B01.P01.C08.S11.SS11.P003}

The concept is relevant to the Philosophy of
Viability because it draws attention to the
conditions under which entities become
manifest within a horizon. It also introduces
the possibility that manifestation is selective
and accompanied by concealment. This is a
valuable comparative problem, but it does
not establish that Heidegger's clearing is
identical with the concept of Manifestation
developed in Part I.
\cite{WrathallHeideggerSEP2025}

% Structural ID: VS.B01.P01.C08.S11.SS11.P004
\phantomsection\label{para:VS.B01.P01.C08.S11.SS11.P004}

A particularly important boundary concerns
the relation between \textit{Lichtung} and
Ontological Light. The German term is usually
translated as ``clearing'' and belongs to
Heidegger's account of disclosure. It must
not be imported as a ready-made definition
of Light in Part III. The ontological status
of Light remains a distinct problem within
Book I, and any proposed relation between
these concepts would require independent
argument.


% =============================================================================
% 12. THE LATER HEIDEGGER: HISTORY, TECHNOLOGY, AND DWELLING
% =============================================================================

\subsection{The Later Heidegger: History, Technology, and Dwelling}
\label{subsec:ch08-heidegger-later}

% Structural ID: VS.B01.P01.C08.S11.SS12.P001
\phantomsection\label{para:VS.B01.P01.C08.S11.SS12.P001}

In the later work, Heidegger increasingly examines
the history of Being and the changing horizons
within which entities are disclosed. The history
of Being is not simply a chronology of philosophical
opinions. It concerns transformations in the
prevailing ways in which entities are understood,
ordered, and encountered. These historical
configurations shape what can appear as real,
significant, available, or intelligible within
a particular epoch.
\cite{WrathallHeideggerSEP2025}

% Structural ID: VS.B01.P01.C08.S11.SS12.P002
\phantomsection\label{para:VS.B01.P01.C08.S11.SS12.P002}

The critique of modern technology is directed
at a mode of revealing, not merely at machines
or technical devices. Heidegger uses
\textit{Gestell} to name the ordering framework
through which entities are disclosed as
standing-reserve (\textit{Bestand}): resources
available for ordering, calculation, storage,
and use. The danger lies in the possibility
that this mode of revealing becomes dominant
and conceals other ways in which entities
may be encountered.
\cite{WrathallHeideggerSEP2025}

% Structural ID: VS.B01.P01.C08.S11.SS12.P003
\phantomsection\label{para:VS.B01.P01.C08.S11.SS12.P003}

This critique should not be reduced to a
general rejection of technology. Heidegger's
question concerns the ontological style
within which technological ordering becomes
the governing mode of disclosure. A device
may be useful without every encounter with
the world being governed by enframing.
The distinction is between the existence
of technical equipment and the dominance
of a particular mode of revealing.
\cite{WrathallHeideggerSEP2025}

% Structural ID: VS.B01.P01.C08.S11.SS12.P004
\phantomsection\label{para:VS.B01.P01.C08.S11.SS12.P004}

The later account of dwelling and place also
bears on the question of relational constitution.
In ``Building Dwelling Thinking,'' Heidegger
examines place as a gathering that cannot be
adequately reduced to a point within a
geometrical coordinate system. A bridge, for
example, can gather a region into a place
through the relations it establishes among
the river, banks, paths, and human dwelling.
This account is relevant to the earlier
manuscript's discussion of spatial identity,
but the gathering of a place must not be
equated with the identity of that place.
\footnote{Martin Heidegger, ``Building Dwelling
Thinking,'' in \textit{Poetry, Language, Thought},
trans. Albert Hofstadter (New York: Harper \&
Row, 1971). \cite{HeideggerPoetryLanguageThought1971}}

% Structural ID: VS.B01.P01.C08.S11.SS12.P005
\phantomsection\label{para:VS.B01.P01.C08.S11.SS12.P005}

The later work also raises historical and
political questions that cannot be ignored
in a responsible examination. Heidegger's
support for National Socialism, his antisemitic
statements, and his later efforts to obscure
the character of his political involvement
are documented matters of philosophical
and historical significance. The relation
between these commitments and his accounts
of historical disclosure and technology
remains a subject of serious dispute. A
philosophical assessment must neither
erase this history nor assume, without
argument, that every concept in his work
is thereby settled or refuted.
\footnote{For the documented political history
and the debate over its philosophical
significance, see Wrathall, ``Martin Heidegger,''
section 3. \cite{WrathallHeideggerSEP2025}}


% =============================================================================
% 13. IDENTITY, DETERMINATION, AND MANIFESTATION
% =============================================================================

\subsection{Identity, Determination, and Manifestation}
\label{subsec:ch08-heidegger-comparison}

% Structural ID: VS.B01.P01.C08.S11.SS13.P001
\phantomsection\label{para:VS.B01.P01.C08.S11.SS13.P001}

Heidegger's analysis contributes to the present
ontology by demonstrating that the intelligibility
of entities cannot always be explained through
isolated properties. Worldhood, practical
involvement, sociality, and temporality disclose
relational contexts within which entities can
be encountered and understood. This gives
substantive weight to the investigation of
Determination and Manifestation as problems
that cannot be treated as detached catalogues
of features.

% Structural ID: VS.B01.P01.C08.S11.SS13.P002
\phantomsection\label{para:VS.B01.P01.C08.S11.SS13.P002}

Yet Heidegger's account does not provide a
direct definition of Identity. His analysis
of Dasein's selfhood concerns the existential
way in which Dasein owns or evades its
possibilities. The constancy of the self
through resoluteness concerns an existential
mode of selfhood, not a complete account
of the ontological identity of every existent.
The distinction between selfhood and Identity
must therefore remain explicit.

% Structural ID: VS.B01.P01.C08.S11.SS13.P003
\phantomsection\label{para:VS.B01.P01.C08.S11.SS13.P003}

Heidegger also places pressure on the relation
between identity and manifestation. If entities
are encountered within horizons of disclosure,
then the manner in which they appear is not
an irrelevant surface added to an otherwise
fully intelligible object. But the dependence
of intelligibility on disclosure does not by
itself prove that the Being or Identity of an
entity is produced by its disclosure. The
relation between ontological constitution
and the conditions of manifestation remains
to be argued.
\cite{WrathallHeideggerSEP2025}

% Structural ID: VS.B01.P01.C08.S11.SS13.P004
\phantomsection\label{para:VS.B01.P01.C08.S11.SS13.P004}

The distinction between Being and beings
also exposes a question about the scope of
the present ontology. Book I examines Being
and Identity as analytically distinguishable
but ontologically inseparable. Heidegger's
question concerns the meaning of Being in
a way that is not exhausted by the analysis
of identity. The comparison therefore suggests
that an ontology of Identity must remain
attentive to the question of Being without
assuming that Heidegger's answer can simply
be adopted as its own.

% Structural ID: VS.B01.P01.C08.S11.SS13.P005
\phantomsection\label{para:VS.B01.P01.C08.S11.SS13.P005}

The relation between the two projects is
thus one of critical proximity without
equivalence. Both resist reducing ontology
to a list of externally related properties.
Both give serious attention to relational
and structural conditions. But Heidegger's
fundamental ontology asks after the meaning
of Being, whereas Book I investigates the
ontological status of Identity and its
relation to Being and Manifestation. The
overlap is real, but the questions and
explanatory tasks remain distinct.


% =============================================================================
% 14. THE STRONGEST CHALLENGE AND BILATERAL EVALUATION
% =============================================================================

\subsection{The Strongest Challenge and Bilateral Evaluation}
\label{subsec:ch08-heidegger-evaluation}

% Structural ID: VS.B01.P01.C08.S11.SS14.P001
\phantomsection\label{para:VS.B01.P01.C08.S11.SS14.P001}

The strongest Heideggerian challenge begins
with the question of whether an ontology
can adequately describe beings while leaving
the meaning of Being insufficiently clarified.
If the concepts of Identity, Determination,
and Manifestation are applied to entities,
the ontology must explain the sense in which
those entities are said to be. Otherwise,
its conceptual distinctions may remain
operative within an unexamined horizon
of Being. This challenge concerns the
foundations of the inquiry, not merely
its terminology.
\cite{HeideggerBeingTime}

% Structural ID: VS.B01.P01.C08.S11.SS14.P002
\phantomsection\label{para:VS.B01.P01.C08.S11.SS14.P002}

A second challenge concerns relational
constitution. Heidegger can ask whether
the ontology of Viability gives sufficient
weight to the practical, social, and
historical contexts within which entities
are encountered. If the account of an
entity begins with a set of determinations
and only later adds relations, it may
overlook the possibility that some
determinations are intelligible only
within a world of involvement. This
objection requires a substantive response,
not a merely terminological distinction.

% Structural ID: VS.B01.P01.C08.S11.SS14.P003
\phantomsection\label{para:VS.B01.P01.C08.S11.SS14.P003}

The table below summarizes the principal
Heideggerian distinctions, their explanatory
contributions, and the limits of the
conclusions that can be transferred to
the ontology developed in Book I.

% =============================================================================
% Project : THE PHILOSOPHY OF VIABILITY
% Book    : Book I — Ontology of Viability
% Book ID : VS.B01
% Part    : Part I — Identity and Manifestation
% Part ID : VS.B01.P01
% Chapter : Chapter 08 — Comparative Philosophical Examination
% Chapter ID : VS.B01.P01.C08
%
% File :
%   4_ASSETS/tables/
%   VS.B01.P01.C08.S11.SS14.P003.TL01.heidegger_claims_and_ontological_limits.tex
%
% Section : Section 11 — Martin Heidegger
%
% Asset :
%   Table TL01 — Heidegger's Claims and Ontological Limits
%
% Structural Identifier :
%   VS.B01.P01.C08.S11.SS14.P003.TL01
%
% Asset Type : TL
%
% Purpose :
%   To distinguish Heidegger's principal concepts, their
%   explanatory function, and the limits of their transfer
%   into the Ontology of Viability.
%
% Intellectual Benefit :
%   Provides a compact comparative reference and prevents
%   the conflation of Being, beings, Dasein, Identity,
%   worldhood, disclosure, and temporality.
%
% Aesthetic / Experiential Benefit :
%   Summarizes the extended discussion in a readable
%   comparative matrix.
%
% Representational Limit :
%   This table is an analytical summary. It does not claim
%   that Heidegger's work is internally uncontested.
%
% Dependencies :
%   longtable, booktabs, array, project column types.
%
% =============================================================================

\small
\renewcommand{\arraystretch}{1.18}
\setlength{\tabcolsep}{4pt}

\begin{longtable}{
    >{\raggedright\arraybackslash}p{2.15cm}
    >{\raggedright\arraybackslash}p{0.25\textwidth}
    >{\raggedright\arraybackslash}p{0.25\textwidth}
    >{\raggedright\arraybackslash}p{0.25\textwidth}
}

\caption{Heidegger's Principal Concepts and Their Ontological Limits}
\label{tab:VS.B01.P01.C08.S11.SS14.P003.TL01} \\

\toprule
\textbf{Concept}
&
\textbf{Heidegger's account}
&
\textbf{Explanatory contribution}
&
\textbf{Limit for Book I}
\\
\midrule
\endfirsthead

\toprule
\textbf{Concept}
&
\textbf{Heidegger's account}
&
\textbf{Explanatory contribution}
&
\textbf{Limit for Book I}
\\
\midrule
\endhead

\midrule
\multicolumn{4}{r}{\footnotesize\textit{Continued on next page}}\\
\endfoot

\bottomrule
\endlastfoot

Ontological difference
&
Distinction between Being and beings.
&
Prevents Being from being treated as one
more entity or highest object.
&
Not equivalent to the distinction between
Being and Identity.
\\

\addlinespace

Dasein
&
The entity that we ourselves are, whose
Being is an issue for it.
&
Clarifies the existential conditions of
the question of Being.
&
Does not provide a universal definition
of every existent.
\\

\addlinespace

Being-in-the-world
&
A unitary structure of Dasein's existence.
&
Shows the significance of worldhood,
practical involvement, and relational context.
&
Does not establish that all beings are
constituted by human practical involvement.
\\

\addlinespace

Readiness-to-hand
&
Equipment encountered in practical use
and involvement.
&
Reveals referential contexts obscured
by detached inspection.
&
A mode of encounter is not identical
with the Identity of an entity.
\\

\addlinespace

Being-with and das Man
&
Sociality and anonymous public norms
are existential structures of everydayness.
&
Shows that self-understanding is socially
and practically situated.
&
Social roles and public norms do not
by themselves determine the number
of identities.
\\

\addlinespace

Care
&
The structural whole of Dasein:
projection, thrownness, and involvement.
&
Connects facticity, possibility, and
practical engagement.
&
Not a definition of Identity or a
universal structure of all beings.
\\

\addlinespace

Authenticity
&
Taking over one's own finite possibilities.
&
Clarifies an existential mode of selfhood.
&
Not a proof of numerical identity,
substantial selfhood, or moral superiority.
\\

\addlinespace

Temporality
&
Ecstatic unity of future, having-been,
and present.
&
Shows the temporal structure of Dasein
and the horizon of its possibilities.
&
Not identical with continuity, persistence,
or identity through transformation.
\\

\addlinespace

Lichtung
&
The open region of disclosure and concealment.
&
Clarifies the relation between unconcealment
and the intelligibility of entities.
&
Not identical with Manifestation or
Ontological Light in Book I.
\\

\addlinespace

Gestell
&
The technological mode of revealing
that orders entities as standing-reserve.
&
Critiques the dominance of one mode
of disclosure.
&
Not a general rejection of technology
or a complete ontology of material beings.
\\

\addlinespace

Dwelling and place
&
Place is understood through gathering
and dwelling, not merely coordinates.
&
Challenges a purely geometrical account
of spatial situatedness.
&
Gathering does not by itself establish
the Identity of a place.
\\

\end{longtable}

\normalsize

% =============================================================================
% END OF TABLE TL01
% =============================================================================

% Structural ID: VS.B01.P01.C08.S11.SS14.P004
\phantomsection\label{para:VS.B01.P01.C08.S11.SS14.P004}

The response of Book I must begin by
recognizing the force of the question
of Being and the significance of
relational contexts. The analytical
distinction between Being and Identity
does not license the treatment of
Identity as an independent object
detached from Being. Nor does the
ontological inseparability of Being
and Identity imply that their
analytical distinction is dispensable.
The three principles remain separate
and must be defended through their
own arguments.

% Structural ID: VS.B01.P01.C08.S11.SS14.P005
\phantomsection\label{para:VS.B01.P01.C08.S11.SS14.P005}

The comparison has not established
that Heidegger's existential analytic
is the foundation of all ontology.
It has established a serious challenge:
the ontology of Identity must clarify
how its account of Being relates to
the conditions under which entities
are disclosed and understood. The
scope of this challenge remains
larger than the present comparison.
It is an argued open problem, not
a resolved defect and not a new
foundational commitment.


% =============================================================================
% 15. TRANSITION TO AL-FARABI
% =============================================================================

\subsection{Transition to Al-Farabi}
\label{subsec:ch08-heidegger-transition}

% Structural ID: VS.B01.P01.C08.S11.SS15.P001
\phantomsection\label{para:VS.B01.P01.C08.S11.SS15.P001}

Heidegger reopens the question of Being
by distinguishing Being from beings and
by examining the existential structures
through which Dasein encounters a world.
His analysis reveals the significance of
practical involvement, social relations,
finitude, and temporality for the
intelligibility of entities. These
contributions place genuine pressure
on the ontology of Identity without
providing a ready-made definition
of Identity itself.

% Structural ID: VS.B01.P01.C08.S11.SS15.P002
\phantomsection\label{para:VS.B01.P01.C08.S11.SS15.P002}

The next comparison turns to Al-Farabi.
The transition is not a movement from
one equivalent vocabulary to another.
It is a movement between distinct
philosophical projects: Heidegger's
question of Being and the existential
analytic, and Al-Farabi's metaphysical
account of existence, causality,
intellect, and the ordered structure
of reality. The terms of the latter
must be reconstructed from their own
texts and conceptual setting.

% Structural ID: VS.B01.P01.C08.S11.SS15.P003
\phantomsection\label{para:VS.B01.P01.C08.S11.SS15.P003}

The question carried forward is whether
an ontology can clarify the structure
and relations of beings without
presupposing an answer to the question
of Being, and how different metaphysical
traditions understand the relation
between the existence of entities,
their determinations, and the order
within which they are intelligible.
The following section examines
Al-Farabi's response on its own terms.

% =============================================================================
% END OF SECTION 11 OF CHAPTER 08
% =============================================================================